\documentclass[12pt,a4paper]{article}
\usepackage[scheme=plain]{ctex}
\usepackage[margin=2.5cm]{geometry}
\usepackage{indentfirst}
\usepackage{setspace}
\usepackage{fontspec}
\setCJKmainfont[BoldFont=SimHei,ItalicFont=KaiTi]{SimSun}
\newCJKfontfamily\kaifont{KaiTi}[AutoFakeBold]
\usepackage{newpxtext,newpxmath}

\usepackage{lastpage}





\pagestyle{plain}

\usepackage[nobottomtitles*]{titlesec}
\titleformat{\section}{\Large\bfseries}{}{0em}{ }
\titlespacing*{\section}{0pt}{3ex}{1.5ex}

\usepackage{xcolor}
\usepackage{needspace}

\setlength{\parindent}{2em}
\setlength{\parskip}{0.4em}

\doublespacing

\title{\bfseries Day 2\\[0.3em]泊入:机遇、抉择与身份}
\author{}
\date{}

\begin{document}
\maketitle

\vspace{-2em}

\vspace{1em}

\section{前哨的第二封信}
\vspace{-0.3em}{\color{gray}\small\kaifont\noindent K 又在村庄度过了两个白天和夜晚。两个可疑的"老助手"不请自来；城堡的第一封公函经由信使送到手上；在贵宾楼的酒吧里，他成了第十办公厅主任的情人的情人。K 在一步步接近城堡，也在一步步被城堡吸纳。}\vspace{0.5em}
\vspace{0.3em}

\noindent\rule{\textwidth}{0.4pt}

\vspace{0.3em}



同志们：



事情的发展比我预期的要快得多，也复杂得多。



先说我获得的东西。两个助手出现了，自称城堡派来的，对土地测量一无所知。我是怎么区分他们的？我用同一个名字叫他们俩。给他们任何任务，都共同承担全部责任。要么他们拆穿彼此的谎言，要么在这件事上成为同一个人。哪一种结果对我都有利。



接着信使来了。巴纳巴斯，一个举止和眼神都与旅馆里那些农民完全不同的人。他把第十办公厅主任的信交到我手上。我独自在阁楼房间里反复读了多遍。信中的措辞很微妙，既把我当成可以自由决定是否供职的人，又步步为营地把我塞进"下属"的位置。如果我接受这个位置，我猜我的行动空间会比现在大得多。不接受呢？信中没写后果。不写比写更重。



然后是弗丽达。克拉姆的情人。我透过门孔第一次看到了克拉姆本人，一个富态的、戴着夹鼻眼镜的老爷，坐在书桌后面，桌上没有文件。一束刺眼的灯光打在他身上。随后，我对弗丽达说：离开克拉姆，做我的情人。我说出口了。她答应了。



但我现在坐在这里写这些字的时候，有一个问题挥之不去。我对弗丽达说那些话，有多少是因为她这个人，有多少是因为她是克拉姆的情人？我当时以为我能把这两件事分开，现在不这么确定了。或者说，在这片田野里，情感本身就是策略的一部分，区分它们是徒劳的。



K

\newpage
\section{老助手？老助手。}
\vspace{-0.3em}{\color{gray}\small\kaifont\noindent K 回到酒店，两个自称城堡派来的"老助手"已经在等着他。他们没有仪器，不懂土地测量，而 K 照单全收。}\vspace{0.5em}
\vspace{0.3em}

\noindent\rule{\textwidth}{0.4pt}

\vspace{0.3em}



当他们快到酒店——K还记得那近处的一个拐弯——时，他大为惊异：天竟然完全黑了。他离开这里已经那么长时间了吗？可是他估算一下，大约才不过一两个小时而已，他是早上离开酒店的，现在一直还没有感到有吃饭的欲望，不久之前又还是大白天，现在可就突然漆黑一片了。"多短的白天哟！多短的白天哟！"他自言自语着，从雪橇上就势滑下来，向酒店走去。

令他十分欣慰的是，酒店老板已经站在门前的台阶上迎接他，手里还举着风灯为他照路。K走了几步又匆匆记起车夫而停了一下，黑暗中什么地方有人连连咳嗽，那就是他了。先不管他了吧，反正很快就会再见到他的。等K上了台阶，走到毕恭毕敬地向他敬礼的店老板面前时，才发现他左右两侧各站着一个人。他从老板手里接过了风灯，就着光亮细看这两个人：原来这正是他在路上遇到过的那两个被称为阿图尔和耶里米亚的男子。现在他们向他举手行军礼。这使他回想起他服役的日子，想起那些快活的岁月而笑了。"你们是什么人？"他问，同时一会儿看看这个，一会儿看看那个。"您的助手，"他们回答。"他们是您的助手，"店老板低声证实道。"怎么着？"K问，"你们就是我那两个老助手，就是我让你们随后来，我在这里等着你们的那两个人吗？"两人点头称是。"很好，"片刻之后K说道，"你们来了很好。"——"不过，"又过了片刻，K说，"你们迟到的太多了，你们太拖拉了。"——"路很远呵，"一个说。"路是很远，"K重复说，"可是你们从城堡来时我碰上你们了。"——"对，"他们说，没有更多的解释。"仪器你们都放在哪儿了？"K问道。"我们没有仪器，"他们说。"就是我交给你们保管的那些仪器，"K说。"我们没有，"他们又重复一遍。"唉，你们是怎么搞的！"K说，"你们懂不懂什么叫土地测量？"——"不懂，"他们说。"如果你们是我的老助手，那么你们就必须懂这个，"K说。他们不言语了。"好，那么来吧，"K说着便推着两人进了酒店。

之后，他们三个几乎是一声不响地坐在店堂里一张小桌旁，啤酒摆在桌上，K坐在中间，两个助手一左一右。除他们外，店里现在就只有另一张桌旁坐着几个农民，与头天晚上差不多。"同你们相处可难了，"K说，一边来回比较两人的脸庞，像他多次做过的那样，"我到底怎么分清你们两个呢？你们只有名字不同，除此以外，你们长得太像了，没有一点区别，就跟"——他停顿了一下，然后便不由自主地说下去——"除此以外你们长得就跟两条蛇一样相像。"他们微笑了。"除了名字以外人们也是很容易区别我们的，"他们分辩说。"这我相信，"K说，"我亲眼看见了，可是我呢，我只能用我的眼睛，而光用眼睛是无法分清你们两个谁是谁的。因此，我将把你们当成一个人看待，管你们都叫阿图尔，你们不是有一个叫这个名字吗？是不是你？"K问其中一个。"不是，"这人答道，"我叫耶里米亚。"——"这没什么关系，"K说，"我将管你们两个叫阿图尔。要是我说派阿图尔去哪里，那么你们两人都去，要是我让阿图尔做某件事，那么你们就两人都去做这件事。这对我虽然有一个很大的不便，就是我不能分别使用你们，然而却也有一大优点，就是对我交给你们做的任何工作，你们都共同担负全部责任。你们自己内部怎样分工我不管，但你们不能互相推诿，对我来说你们只是一个人。"他们考虑了一会儿，然后说道："这样我们觉得太别扭了。"——"怎么能不别扭呢，"K说，"这对你们当然一定会很别扭，可就得这样定下来。"

这一阵子以来，K早看见一个农民轻手轻脚地在他们桌旁转悠，探头探脑，欲言又止，现在他终于下决心走到一个助手身后，想对他说几句悄悄话。"对不起，"K说，同时一捶桌子，站了起来，"这两位是我的助手，我们现在正在谈正事。谁都没有权利打扰我们。"——"哦，请便，哦，请便，"那农民战战兢兢地说，边说边往后退着回到和他一起的那帮人那里去了。"这一点你们两个务必注意，"K重新坐下后说道。"没有我的许可，你们不得同任何人谈话。我在这里是个外乡人。既然你们是我的老助手，那你们也就是外乡来的。因此，我们三个外乡人要齐心，伸出你们的手，让我们拉拉手吧。"两人忙不迭地向K伸出手来。"你们两个也握一握手吧，"他说，"对我要有令则行。我现在要去睡了，我劝你们也去睡一觉。今天我们耽误了一个工作日，明天必须很早开始工作。你们要设法找一辆雪橇来，我们好去城堡，六点钟必须一切准备停当在这门口待命。"——"好的，"其中一人说。但另一个却抢白说："你说‘好的‘，可你明明知道这办不到。"——"不要说了，"K说，"难道你们现在就要一个跟一个闹分裂？"但是，现在第一个助手也否认了："他说得对，没有许可任何外人是不得进入城堡的。"——"要到哪里去申请许可呢？"——"我不知道，也许是找主事吧。"——"那我们就打电话到那儿去申请，你们两个听着：马上打电话给主事！"两人立即跑步来到电话机旁，接通了主事——瞧他们那争先恐后的劲头儿！从表面看，他们真是唯命是从、唯唯诺诺到可笑的地步——问道，K可不可以同他们一起上城堡。电话里厉声回答："不行！"就连K坐在他桌旁也听见了。然而答话比这还要多一点，叫做："明天不行，任何时候都不行。"——"我自己来打，"K说着站了起来。到目前为止，除了那个农民引起的干扰外，很少有人注意K和他的两个助手，然而现在他这最后一句话却引起了众人的注意。所有的人一下子跟他一齐站了起来，然后，虽然老板力图把这些人推开，他们仍在电话机旁围成一个半圆形的圈子将K包围了。他们议论纷纷，占主导地位的看法是：K将得不到任何回答。K不得不请求他们安静下来，说他并没有要求他们把他们的意见讲给他听。

听筒里传来K以往打电话时从未听到过的嗡嗡声。听起来，就像从一大片乱哄哄的孩子吵嚷声中——可这声音又不是真正的嗡嗡声，而是从远方，从很远很远的地方传来的歌唱声，就像是从这一片嗡嗡声中神奇而不可思议地逐渐幻化出一个单一的、很高的强音，它猛烈撞击着他的耳鼓，仿佛它强烈要求深深钻入人体内部而不只是接触一下那可怜的听觉器官似的。K全神贯注地等候，一句话不说，左臂支在摆电话机的小桌上，就这样聚精会神地恭候着。

他不知过了多长时间；只知道后来老板拉扯他的上衣，告诉他有一个信差来找他。"走开！"K忍不住大叫一声，也许这吼声传进听筒里去了吧，因为那边马上有人接电话了。于是他听见了下面的话："我是奥斯瓦尔德，那边是谁呀？"听筒里叫道，这是一个严厉而高傲的声音，K觉得有点小小的语音缺陷，那声音力图超越自身，用更加咄咄逼人的凌厉气势来弥补这个缺陷。K在犹豫，没有立即报上自己的名字，在电话面前他是赤手空拳，毫无防御能力的，对方可以厉声呵斥他，可以把听筒撂下，这样一来K就无异于自己把一条也许是相当重要的通路堵死了。K的犹豫使那边的男子不耐烦起来："打电话的谁呀？"重复喝问之声接着又补充道："请你们那边不要老打电话来好不好？刚刚才来过一个电话嘛！"K没有理睬这句话，他突然灵机一动，报上自己身份说："我是土地测量员先生的助手。""什么助手？什么先生？什么土地测量员？"K想起了昨天的电话。"请您问问弗里茨吧，"他简短地说。使他吃惊的是这话居然奏效了，比这更使他惊讶的是那边办事口径完全一致。答话是："我知道了。那个永远扯不清的土地测量员。是了，是了。还有什么？哪个助手？""约瑟夫，"K说。他对那些农民在他身后嘀嘀咕咕有些不快；显然他们不同意他没有报上真实身份这一做法，但K没有功夫去考虑他们，因为应付电话就够使他费神的了。"约瑟夫？"那边反问道。"两位助手不是叫"——说到这里那边稍稍停了一下，显然是在让谁告诉名字——"阿图尔和耶里米亚吗？""他们是新来的助手，"K说。"不，他们是老助手。"——"他们是新的助手，我才是老助手，是今天比土地测量员先生晚一步后到这里的。"——"不对！"那边大叫起来了。"那么我是谁呢？"K问道，一直保持着冷静。这样过了一会儿，听筒里又响起同一个声音，带着同样的语音缺陷，可同时又的确像是另一个更低沉、更令人肃然起敬的声音："你是老助手。"

K过于专注地琢磨那声音的韵味，以至几乎没有听见对方的问话："你想干什么？"〔2〕现在他真恨不得立刻放下听筒。他已不再期望能从这次谈话中得到点什么。他只是硬着头皮匆匆再问了一句："什么时候我的主人可以到城堡来？"——"什么时候都不行，"就是回答。"知道了，"K说着便挂上了听筒。
\par\vspace{0.3em}{\footnotesize\color{gray}（第二章）}
\newpage
\section{一封危险的公函}
\vspace{-0.3em}{\color{gray}\small\kaifont\noindent 信使巴纳巴斯，一个看似超脱于村庄的年轻人，递上了第十办公厅主任署名的公函。信中的措辞模棱两可：K 是做微末的下属，还是主任身边的透明人？}\vspace{0.5em}
\vspace{0.3em}

\noindent\rule{\textwidth}{0.4pt}

\vspace{0.3em}



这时他身后那些农民早已挤到他跟前来了。两个助手一边不断斜眼看K，一边不住地挡住农民不让他们靠近。然而他们只不过是做做样子罢了。农民们实际上也已经对谈话结果感到满意而逐渐后退了。正在这时，从他们后面，一个人大步流星，把农民向两边分开走了过来，他向K鞠了一躬，递交给他一封信。K手里拿着信，两眼注视着来人，觉得这人眼下比别的事情更加重要。他和两个助手有许多相似之处，同他们一样身材瘦长，一样穿着紧身衣，也同他们一样敏捷灵巧，可又与他们完全不同。要是我K的助手是这个人多好！这人某些地方使K想起在盖尔斯泰克那里见过的那个怀抱婴儿的女人。他的着装几乎全是白色，那套衣服大概不是丝绸的，它是很普通的冬衣，但却具有丝绸衣服的柔和和庄重。他的脸明亮而开朗，眼睛非常大。他的微笑有出奇的感染力；他以手拂面，那神气似乎想赶走这微笑，然而却不成功。"你是谁？"K问道。"我叫巴纳巴斯，"他说。"我是个信使。"说话时他那两片嘴唇一张一闭，既显出阳刚之气，又饱含柔和之美。"你喜欢这里吗？"K问道，同时指了指那伙农民，到现在他一直还保持着对他们的兴趣，这些人有着一张张简直就是受苦受难的脸——他们的颅骨似乎被人打扁了，面部线条则是挨打后的痛苦表情勾勒出来的——和厚大的翻嘴唇，他们张开大嘴站在一旁观看，但同时又没有观看，因为有时他们的目光游移不定，茫然若失地盯住某件无关紧要的东西瞅一阵，然后才又返回到他和巴纳巴斯身上来；除这些农民外，K又指了指他的两个助手，这两人现在紧紧地手拉着手，脸贴着脸，面带微笑，无法确定这笑是卑逊的呢还是轻蔑的，他把他们全都指给他看，仿佛在介绍特殊境遇强加于他的一批随从，并期望着——这期望里有一种一见如故的心情，K很希望同这个人谈谈心里话——巴纳巴斯时时把他和这些人区别开来。但是巴纳巴斯根本就不回答——当然不含任何恶意，这是看得出来的——他的问题，恰似一个训练有素的仆人在听到主人仅仅表面上针对他而说的话时露出的表情那样，顺从地默默听着，只是遵循所提问题的意思左顾右盼，然后举手向农民中的熟人打招呼，又同两个助手交谈了几句，言谈、举止、神态是自由而自主的，一点没把自己同他们混同起来。K转眼——他问话碰了个软钉子，但并不觉难堪——看看手中的信并拆开了它。这封信的全文是："非常尊敬的先生！如您所知，您已被聘任前来为大人供职。您的直接上司是村长，他还将告知您有关您的工作及薪俸的一切细节，而您也将负责向他汇报工作。尽管如此，我本人还是要在百忙中兼顾您的。递送此信的巴纳巴斯将不时向您询问以了解您的愿望，然后及时禀告我。您将发现我会尽一切可能为您提供方便。使我的部下心满意足，实为我所期盼。"签署的名字很潦草，看不清楚，但在签署旁边印有："第十办公厅主任"字样。"你等一等！"K对正向自己鞠躬的巴纳巴斯说，然后他叫老板要一个房间用用，说他想一个人再好好看看这封信。同时他又想到，不论自己对巴纳巴斯有多少好感，他毕竟只是个信差，于是吩咐拿些啤酒给他。K留意观察他对此的反应，看到他显然很高兴地接受了这一赏赐，立刻喝起来。然后K便随老板而去。在酒店这所小房子里，只能做到为K准备一间阁楼斗室，而即使这一点也不那么容易，因为必须把原来在这房间里睡觉的两个女仆安排到别处去住。其实，也只不过是让她们出去而已，房间本身看来原封未动：唯一的一张床上，没有床单、枕套等物，只有几块垫子和一床粗毛毯子，还保持着昨夜使用过以后的原状。墙上挂着几张圣徒画像和几张士兵照片。屋子里甚至没有开窗换换空气，显然是希望新来的客人不要在这里长住，所以不做任何能吸引他留下的事。但K挺能凑合，怎么都行，他披上毯子，坐到桌前，开始借着一支蜡烛的光亮再次读起那封信来。

这封信的内容前后并不一致，有几处像在同一个享有充分自由的人谈话，他自己的意愿受到尊重，如信头，还有涉及他的愿望的那句话都是如此。然而信中也有一些地方或隐晦或明显地把他当成一个小小的、从办公厅主任坐椅上看去小得几乎看不见的下属，主任必须作出很大的努力，才能"在百忙中兼顾"他一下，他的上司只是村长，K甚至还有向这位村长汇报工作的责任，他的唯一同事恐怕就只有那个村警了。这无疑都是一些矛盾，它们过于明显，使人觉得一定是有意这样措词的。K几乎不敢在这样的官府面前不知天高地厚地认为这里可能有犹豫不决的因素在起作用，相反，他觉得信中是明白无误地给他提出了两种可能的选择，完全让他自己决定如何对待信中作出的安排：是想做一名同城堡有着说起来倒是很光彩、然而却只是一种表面上的联系的农村工作人员呢，还是做一名仅仅是表面上的农村工作人员，而事实上却由巴纳巴斯带的信息来决定他和城堡的全部工作关系？K在选择上毫不犹豫，而且即便他没有到这里之后取得的那些经验，他也是不会犹豫的。只有作为一名农村工作人员，离城堡老爷们愈远愈好，他才有可能在城堡达到一定的目的。村里这些人，现时对他还抱着很大的不信任，而一旦他不说成为他们的朋友，就算只成为同他们一样的村民，他们就会开口说话，而一旦他同盖尔斯泰克或拉塞曼没有任何区别——这一点必须迅速做到，这是关键的一步，那么，在他面前所有的道路就一定会一下子畅通无阻，而如果他仅仅寄希望于山上那些老爷们及其恩典，他就不仅会永远被阻挡在这些大路之外，而且恐怕连这些路的影子也永世看不到。当然，存在着一种危险，信中反复强调了这种危险，在表述它时字里行间流露出某种喜悦，似乎它是摆脱不掉的了。这危险就是他的下属地位：供职、上司、工作、薪俸、汇报、部下，这封信中充斥着这些字词，就是在谈到别的问题，比较涉及个人的问题时，也是从这一角度出发的。如果K愿意，那么他马上可以成为这个部下，但那就是极为严肃的事，没有丝毫另就的可能。K知道人家并没有拿要采取任何真正的强迫手段来威胁他，他不怕强迫，尤其在这个地方更不怕，但是他害怕这个使人泄气的环境对他产生的强大压力，害怕对各种失望逐渐习以为常的那样一种习惯势力，害怕无时无刻不在向他袭来的各种潜移默化的影响形成的强大压力，然而无论怎样害怕，他必须冒着这种危险勇敢地投入战斗。信上不是也不讳言，如果真需要靠斗争来一决雌雄，那么K是有胆量先发制人的吗；话说得很含蓄，唯有一颗焦灼不安的心——是焦灼不安的心，而不是深感歉疚的心——才能觉出这里的弦外之音，正是"如您所知"那几个关于他被录用为大人供职的字眼，暗含着这层意思。K只是报了到，但自那时起他便如信中所说，已经知道被录用了！

K从墙上取下一张画，把信挂在那个钉子上；既然他将在这间屋子里住，那么信就理应挂在这里。
\par\vspace{0.3em}{\footnotesize\color{gray}（第二章）}
\newpage
\section{双向奔赴的爱情？}
\vspace{-0.3em}{\color{gray}\small\kaifont\noindent 在巴纳巴斯的姐妹奥尔嘉的带领下，K 在贵宾楼酒吧遇见了第十办公厅主任克拉姆的情人弗丽达。他选择向通往城堡权力核心的最短路径发起冲击。他成功了吗？}\vspace{0.5em}
\vspace{0.3em}

\noindent\rule{\textwidth}{0.4pt}

\vspace{0.3em}



在酒吧——一间很大的、当中完全空着的屋子——里，立着一箱箱啤酒，农民们有的在酒桶边，有的在酒桶上靠墙坐着，但他们的模样与K住的那家酒店里那些农民完全不同。他们衣着整洁一些，都穿一色浅土黄粗布，上衣鼓起，裤子紧身。这些男人全是矮个儿，乍一看长相都差不多，都是扁平脸庞，额头、颧骨、下巴突出，却又有多肉的面颊。他们一概少言寡语，几乎纹丝不动地坐着，只用目光尾随新来者，然而眼神又是迟滞、冷漠的。尽管如此，因为人多，又非常安静，他们仍然引起了K一定的注意。他再次拉起了奥尔嘉的胳臂，算是告诉众人他是怎么来到这里的。这酒吧的一角有一个男人站了起来，这是奥尔嘉的一个熟人，他想到这边来，可是K立刻挎起她的胳臂推着她朝另一个方向走了。这一点除她本人外谁都没觉察到，她微微一笑，斜睨他一眼，算是默许他这个行动了。

斟啤酒的是一个名叫弗丽达的少女。这是个头发金黄、脸颊瘦削、有着一双忧伤眼睛、不引人注目的小个子姑娘，可是她的目光却令人吃惊，那是一种特别高傲的目光。当它落在K身上时，K觉得它似乎已经把所有与他有关的事情统统解决了，他本人现在还一点不知道有这些事情，然而这目光却使他坚信确有其事。K一直不停地从侧面看着弗丽达，就是当她已经在同奥尔嘉说话时也还在打量她。看样子，奥尔嘉和弗丽达可不大像是朋友，她们只是冷冷地交谈了几句。K想帮她们缓和一下气氛，所以突然动问道："您认识克拉姆先生吗？"奥尔嘉忍不住笑起来。"你笑什么？"K生气地问。"我并没有笑呵，"嘴上这么说，实际上仍不停地笑。"奥尔嘉这姑娘还真够孩子气的，"K说，同时靠着柜台把身子使劲向下弯，以便再次把弗丽达的注意力吸引到自己身上来。但她仍低垂眼皮，轻声说道："您想见克拉姆先生吗？"K说请求一见。她便指着一扇正好就在她自己左手边的门。"这门上有个小孔，您可以从这儿往里面看。"——"那么这儿的这些人呢？"K问道。她撅了撅下嘴唇，用她那异常柔软的手把K拽到门边。把眼睛贴在那个显然是为了窥视而钻出的小孔上，他几乎可以将隔壁房间一览无余。

屋子中央，一张书桌旁有一把舒适的圆形靠椅，椅子上坐着的正是克拉姆先生，他面前挂着一盏亮得刺眼的电灯。这是一位中等身材、颇为富态、看来一定行动不甚方便的老爷。他的脸还算光溜，但面部肌肉却已经由于年龄的分量而有些下垂了。黑色的小胡子长长地伸向两边。一副歪戴着的不断反光的夹鼻眼镜挡住了他的眼睛。如果克拉姆先生完全伏案工作，那么K就只能看见他的侧面；但因为他现在几乎是面对K坐着，所以K可以看清他的整个面庞。克拉姆左肘支在桌上，拿着一支弗吉尼亚雪茄的右手搁在膝盖上。桌上摆着一个啤酒杯；由于桌子的边条较高，K不能看清桌上是否有任何文件之类，但他觉得桌面似乎是空的。为保险起见，他请弗丽达也从小孔中看一看然后告诉他究竟有还是没有。但是，因为她不久前刚到过那屋里，所以便确有把握地向他证实那桌上并没有什么文件。K问弗丽达他现在是否必须走开，可是她说如果他乐意，趴在门上看多久都行。现在K身边只有弗丽达，奥尔嘉呢，K匆匆注意到她已经找到了她的熟人，这时高高坐在一只桶上，脚不停地荡来荡去。"弗丽达，"K小声说，"您同克拉姆老爷一定很熟啰！"——"唔，是的，"她说。"是很熟。"她靠墙站在K身边，并且——这时K才注意到——用手轻轻地抚弄她那件薄薄的肉色开领衬衫，这件衣服同她那可怜的瘦小身躯很不相称。俄顷她说："您不记得奥尔嘉刚才的笑了吗？"——"记得，这个淘气鬼，"K说。"不过，"她带着和解的口吻说，"她笑是有道理的。您刚才不是问我是不是认识克拉姆吗？可我恰好是，"——说到这里她不由自主地挺了挺胸脯，同时她那与谈话内容毫无关系的高傲目光再次越过K的头投向别处，"我恰好是他的情人呵。"——"克拉姆的情人，"K说。她点点头。"那么在您面前，"为不让他们之间的气氛过于严肃，K微笑着说，"我必须肃然起敬了。"——"不光是您，"弗丽达和蔼地说，但并不是受到了他那微笑的感染。K想到一个煞一煞她的傲气的办法，现在使了出来；他问道："您到城堡里去过吗？"可这难不倒她，她立即答道："没去过，可是我在这儿这个酒吧里，这不就足够了吗？"她的虚荣心显然非常强，看来此刻正想在K的身上满足一下这种虚荣心。"当然啦，"K说，"在这个酒吧里，您是很擅长做酒店老板工作的嘛。"——"是这么回事，"她说，"可我一开始时是在‘大桥‘酒店当喂牲口的女佣人呀。"——"难道就用这双这么娇嫩的手喂牲口吗？"K带着疑虑的语气说，连他自己也不清楚他这话究竟只是为了说点好听的呢，还是他真的已经被她俘虏了。她的两只手确实小巧而又娇嫩；但是，管它们叫做干瘪瘦弱、索然无味也未尝不可。"当时谁也没有注意这个，"她说，"就是现在——"K疑惑地看着她。她摇摇头不想说下去。"当然，您有您的秘密，"K说，"您是不会同一个刚认识半小时，还没有机会向您介绍自己情况的人谈论这些秘密的。"可是，这话看来说得很不合适，就好像弗丽达本来处在一种对他很有利的迷迷糊糊的状态，而他竟一下子把她摇醒了。她从自己腰带上挂着的皮包里拿出一小块木头把门上小孔堵住，对K说话时显然在竭力抑制自己，不让他觉察她内心的变化："至于说到您，那么我知道您的底细，您是土地测量员，"然后又补充道："现在我得去干活了，"说完便走到柜台后面她的位置上，时不时有一两个酒客站起来，到柜台前让她斟酒。K还想再同她悄悄谈一次，因此从一个架子上取下一只空杯向她走去。"我只想再说一句，弗丽达小姐，"他说，"从一个喂牲口的女佣人，苦干到当上酒吧女招待，是非常了不起的，需要有超出常人的毅力，可是对于这么个不寻常的人，最终目标是不是就到此为止了呢？这是个荒谬的问题。请别笑我，弗丽达小姐，从您的眼睛里我看到的不完全是您过去的奋斗，而更多的是您将来的奋斗。但这个世界上各种阻碍是很大的，目标越大阻碍也越大，所以，希望争取到一个无权无势但也在奋斗着的小人物的帮助，便不是什么可耻的事情了。也许我们两人可以好好地单独谈谈，别有那么多双眼睛盯着我们说话。"——"我不知道您想干什么，"她说，但这次与她的意愿相反，语调中已听不到她在生活中屡次得胜产生的傲气，而只有数不清的失望引起的凄楚了。"您兴许想把我从克拉姆身边拉走吧，我的老天！"她两手合拢一击掌。"您把我看透了，"K说，显得像是被太多的不信任弄得疲惫不堪，"这一点正是我最隐秘的意图。您最好离开克拉姆，做我的情人。好了，现在我可以走了。奥尔嘉！"K叫道。"我们回家了。"奥尔嘉顺从地就势从桶上滑下，可还不能马上摆脱那帮包围她的朋友。这时弗丽达瞪了K一眼，低声说："我什么时候可以同您谈谈？"——"我可以在这儿过夜吗？"K问。"可以，"弗丽达说。"现在就可以留下吗？"——"您先同奥尔嘉离开这里，我好设法支走眼前这帮人。过一会儿您就可以来了。"——"好吧，"K说，焦急地等着奥尔嘉。
\par\vspace{0.3em}{\footnotesize\color{gray}（第三章）}
\newpage
\section{强势田野与"真问题"}
\vspace{-0.3em}{\color{gray}\small\kaifont\noindent 一部分现代政府研究者以研究强势田野自居，强调对"真问题"的识别能力。存在一套"真问题"掩护机制，以及克服这种掩护机制的一般做法吗？}\vspace{0.5em}
\vspace{0.3em}

\noindent\rule{\textwidth}{0.4pt}

\vspace{0.3em}



（四）"强势田野"的真问题识别困境



"强势田野"意味着研究者极易陷入真问题识别困境，具体表现为三种状态：其一，受信息不对称与自身认知局限，研究者无力穿透田野的表象，难以捕捉具有研究价值的核心议题，陷入无问题的迷茫状态；其二，在表层信息供给与应付式互动的误导下，研究者易将非真实存在的议题误判为核心研究对象，陷入伪问题的认知陷阱；其三，研究者可能聚焦非核心、低学术价值的真实议题，误将边缘问题作为主要研究方向。



无问题是指研究者在"强势田野"中难以识别出具有研究价值的真实问题，陷入"看得见却看不清、接触到却摸不透"的困境。这源于以下两个方面：一方面，"强势田野"存在严格的信息筛选与层级化沟通，研究者往往被限定在公开场景与表层互动中，难以触及核心工作流程与真实运作——田野主体出于风险规避、工作保密或程序规范等考量，仅向研究者展示经过包装的标准化场景，如预先安排的观摩环节、统一口径的访谈回应、公开可查的常规资料，而那些隐藏在流程背后的矛盾冲突、政策执行弹性、主体间利益博弈等关键信息并未向研究者公开，使研究者如同"隔着玻璃看风景"，无法感知田野的真实肌理；另一方面，研究者自身的认知局限与身份隔阂增加了发现问题的难度，"强势田野"往往具备高度专业化的运作规则、独特的行业术语与隐性的行动逻辑，使研究者既难以理解田野中各类现象的深层含义，也无法判断哪些现象背后隐藏着值得探究的问题，进而只能被动接收碎片化、表面化的信息，无法将零散现象串联成有逻辑的问题线索，最终陷入"无问题可究"的迷茫状态。更为严重的是，这可能导致研究者将"无问题"误判为田野的真实状态，或将自身的认知盲区等同于田野的平稳运行，进而得出"该田野不存在值得研究的问题"的错误结论，错失挖掘核心研究议题的机会。



伪问题是指研究者在"强势田野"中识别的问题并非真实存在的问题。这源于以下两个方面：一方面，研究者对"强势田野"理解不深，进入田野如同"刘姥姥进大观园"，一切新鲜事物对其都是"问题"；另一方面，田野主体告知研究者一些老生常谈的信息，如工作的推进源于领导重视、建立领导小组、建章立制、群策群力、实现创新等，这些访谈内容实则属于"放之四海而皆准"的信息。"放之四海而皆准"意味着这些可能与研究者所提出的问题不符的内容很可能被误认为是珍贵的访谈材料。与此同时，田野主体可能提供给研究者诸多文本素材，其多以规章制度、工作总结、工作报告为主，这些材料中不乏有效的田野素材，但研究者可能难以从这些材料中感知真正有用的信息——因为许多材料大同小异，而写材料的逻辑又使文本极为凝练。有鉴于此，一些将"某些任务推进何以成功"的原因归结于上述所说建章立制、群策群力等的观点，很可能并未挖到根源，而由此去探寻某个政策执行的逻辑容易发现伪问题。



边缘问题是指研究者在"强势田野"中识别出的真实存在却不值得着重研究的问题。这源于以下两个方面：一方面，研究者聚焦"强势田野"中的非核心工作，这些工作多是各部门的通用型工作，专业化程度低意味着研究者易于上手，有时研究者还会认为这些"新事物"值得研究；另一方面，田野主体提供给研究者若干亟待解决的问题，其大多是真实存在的问题，如部门运作机制的优化。这些问题多为田野主体站在研究者角度所提出的问题：作为高级知识分子，研究者应去解决实践者解决不了的问题。显然，这些在政府部门普遍存在的问题是真实的，但可能鉴于其微观特性、技术的难解决性以及既有研究已经提供可行的方案而不能提炼为高质量的研究问题。值得注意的是，研究者往往会把边缘问题当作与"强势田野"中主体破冰的"第一桶金"，但这存在风险：如果研究者没有完成，那么田野主体可能怀疑研究者能力并减少互动意愿；如果研究者完成，那么田野主体可能持续交给研究者类似的工作，研究者没有时间进行深入的思考；更有甚者，研究者可能将这些问题当作研究问题，陷入"只见树木，不见森林"的困境。



上述三种情形的核心症结在于研究者未能穿透"强势田野"的表象，触及其实质矛盾与深层逻辑，最终导致具有学术价值的真问题缺失。破解这一困境的关键在于习得并运用地方性知识——在政府研究中，地方性知识有助于研究者把握"强势田野"的运作规则、互动惯例与隐性知识，进而突破信息壁垒、纠正认知偏差，其在问题识别、界定、正当化以及任何可能的解决方案生成过程中都扮演重要角色。
\par\vspace{0.3em}{\footnotesize\color{gray}（凌争，政府研究中的"强势田野"与地方性知识，《行政论坛》2026年第1期，第76-85页）}
\end{document}